% !TeX root = ../main.tex
\section{Результаты эксперимента}

После завершения четырёхминутной стабилизации следящих контуров приёмника в $15{:}27$ местного времени ($12{:}27\text{ UTC}$) были зафиксированы параметры всех наблюдаемых навигационных космических аппаратов. Из $41$ прогнозируемого спутника четырёх основных глобальных навигационных систем контрольный приёмник устойчиво захватил и включил в навигационное решение $31$ аппарат. Сравнительная сводка прогнозного и фактически зарегистрированного количества спутников, а также экстремальные значения отношения несущей к спектральной плотности мощности шума ($C/N_0$) по каждой системе представлены в таблице~\ref{tab:summary_results}.

\begin{table}[H]
    \centering
    \small
    \caption{Сравнение прогноза и натурных наблюдений по системам ГНСС}
    \label{tab:summary_results}
    \begin{tabular}{|l|c|c|c|c|}
        \hline
        Система ГНСС & Прогноз & Наблюдалось & Мин. $C/N_0$, дБ-Гц & Макс. $C/N_0$, дБ-Гц \\
        \hline
        GPS & 8 & 7 & 24,1 & 43,4 \\
        \hline
        ГЛОНАСС & 9 & 7 & 21,9 & 41,1 \\
        \hline
        Galileo & 12 & 9 & 22,4 & 42,6 \\
        \hline
        BeiDou & 12 & 8 & 20,8 & 39,8 \\
        \hline
    \end{tabular}
\end{table}

Детальный перечень захваченных спутников с указанием их системных номеров, углов места (возвышения над горизонтом) и измеренных на контрольном приёмнике уровней $C/N_0$ приведён в таблицах~\ref{tab:gps_details}--\ref{tab:bds_details}. Космические аппараты, находившиеся по расчётам ниже угла маски отсечки приёмника ($9^\circ$) или экранированные элементами городской застройки на малых углах места, не сопровождались.

\begin{table}[H]
    \centering
    \small
    \caption{Параметры наблюдаемых спутников системы GPS (диапазон L1)}
    \label{tab:gps_details}
    \begin{tabular}{|c|c|c|c|}
        \hline
        Номер (PRN) & Угол места & $C/N_0$, дБ-Гц & Статус \\
        \hline
        24 & $78^\circ$ & 43,4 & В решении \\
        \hline
        01 & $53^\circ$ & 38,9 & В решении \\
        \hline
        14 & $34^\circ$ & 34,7 & В решении \\
        \hline
        10 & $27^\circ$ & 31,8 & В решении \\
        \hline
        25 & $18^\circ$ & 28,5 & В решении \\
        \hline
        28 & $15^\circ$ & 26,3 & В решении \\
        \hline
        32 & $11^\circ$ & 24,1 & В решении \\
        \hline
        21 & $6^\circ$ & --- & Не захвачен \\
        \hline
    \end{tabular}
\end{table}

\begin{table}[H]
    \centering
    \small
    \caption{Параметры наблюдаемых спутников системы ГЛОНАСС (диапазон L1)}
    \label{tab:glo_details}
    \begin{tabular}{|c|c|c|c|}
        \hline
        Системный номер & Угол места & $C/N_0$, дБ-Гц & Статус \\
        \hline
        761 & $73^\circ$ & 41,1 & В решении \\
        \hline
        735 & $57^\circ$ & 38,2 & В решении \\
        \hline
        754 & $46^\circ$ & 35,9 & В решении \\
        \hline
        757 & $37^\circ$ & 33,4 & В решении \\
        \hline
        745 & $31^\circ$ & 30,6 & В решении \\
        \hline
        760 & $16^\circ$ & 24,8 & В решении \\
        \hline
        707 & $12^\circ$ & 21,9 & В решении \\
        \hline
        762 & $7^\circ$ & --- & Не захвачен \\
        \hline
        765 & $5^\circ$ & --- & Не захвачен \\
        \hline
    \end{tabular}
\end{table}

\begin{table}[H]
    \centering
    \small
    \caption{Параметры наблюдаемых спутников системы Galileo (диапазон E1)}
    \label{tab:gal_details}
    \begin{tabular}{|c|c|c|c|}
        \hline
        Номер (PRN) & Угол места & $C/N_0$, дБ-Гц & Статус \\
        \hline
        E24 & $81^\circ$ & 42,6 & В решении \\
        \hline
        E04 & $52^\circ$ & 38,1 & В решении \\
        \hline
        E15 & $44^\circ$ & 36,0 & В решении \\
        \hline
        E28 & $41^\circ$ & 34,8 & В решении \\
        \hline
        E23 & $32^\circ$ & 31,9 & В решении \\
        \hline
        E06 & $29^\circ$ & 30,4 & В решении \\
        \hline
        E34 & $23^\circ$ & 27,8 & В решении \\
        \hline
        E19 & $17^\circ$ & 25,1 & В решении \\
        \hline
        E09 & $12^\circ$ & 22,4 & В решении \\
        \hline
        E13 & $7^\circ$ & --- & Не захвачен \\
        \hline
        E21 & $5^\circ$ & --- & Не захвачен \\
        \hline
        E36 & $4^\circ$ & --- & Не захвачен \\
        \hline
    \end{tabular}
\end{table}

\begin{table}[H]
    \centering
    \small
    \caption{Параметры наблюдаемых спутников системы BeiDou (диапазон B1)}
    \label{tab:bds_details}
    \begin{tabular}{|c|c|c|c|}
        \hline
        Номер (PRN) & Угол места & $C/N_0$, дБ-Гц & Статус \\
        \hline
        C08 & $55^\circ$ & 39,8 & В решении \\
        \hline
        C26 & $49^\circ$ & 37,6 & В решении \\
        \hline
        C41 & $38^\circ$ & 34,2 & В решении \\
        \hline
        C31 & $33^\circ$ & 32,1 & В решении \\
        \hline
        C32 & $24^\circ$ & 28,9 & В решении \\
        \hline
        C29 & $21^\circ$ & 26,7 & В решении \\
        \hline
        C33 & $16^\circ$ & 23,5 & В решении \\
        \hline
        C02 & $11^\circ$ & 20,8 & В решении \\
        \hline
        C25 & $8^\circ$ & --- & Не захвачен \\
        \hline
        C35 & $6^\circ$ & --- & Не захвачен \\
        \hline
        C38 & $5^\circ$ & --- & Не захвачен \\
        \hline
        C45 & $4^\circ$ & --- & Не захвачен \\
        \hline
    \end{tabular}
\end{table}

\section{Анализ полученных результатов}

Сравнение данных расчётного прогнозирования и натурных измерений демонстрирует высокую степень соответствия состава рабочего созвездия: фактический захват составил $75{,}6\%$ от общего числа спутников, находящихся выше математического горизонта ($31$ из $41$ НКА). Все космические аппараты, имеющие углы места свыше $10^\circ$, были устойчиво захвачены следящими контурами и включены в навигационное решение для всех четырёх систем ГНСС.

Анализ распределения измеренных значений отношения несущей к спектральной плотности шума подтверждает выраженную зависимость энергетики принимаемого радиосигнала от угла возвышения спутника над горизонтом. Максимальные значения $C/N_0$ в диапазоне $39{,}8\text{--}43{,}4\text{ дБ-Гц}$ зафиксированы для космических аппаратов, находящихся вблизи зенитной области.
